(1) Helmholtz自由能满足 $dF_s=-S_s\,dT+\gamma\,dA$，因此固定 $A$ 时 $S_s=-(\partial F_s/\partial T)_A=-A\gamma'(T)$。若常见液体 $\gamma'<0$，表面熵为正；此处只计算界面超额量，不包含体相熵。
(2) $U_s=F_s+TS_s=A[\gamma(T)-T\gamma'(T)]$。固定面积热容 $C_{s,A}=(\partial U_s/\partial T)_A=-AT\gamma''(T)$；若 $\gamma''<0$，该界面热容正。对温度求导时不可遗漏 $\gamma'$ 的抵消项。
(3) 恒温改变球半径，$dF_s=\gamma dA=8\pi R\gamma\,dR$，体积变化 $dV=4\pi R^2dR$。可逆机械平衡要求 $\Delta p\,dV=dF_s$，故 $\Delta p=2\gamma/R$，即液滴内部压力比外部高的Laplace关系。
(4) 若 $\gamma=\gamma_0-s(T-T_0)$，则 $S_s=As$，$U_s=A(\gamma_0+sT_0)$，$C_{s,A}=0$。需在 $\gamma>0$ 的温区使用此线性模型；若 $s>0$，界面熵为正，且面积改变时仍有非零表面能和机械功。
